The optimization of task offloading in vehicular networks represents a critical challenge for enabling intelligent transportation systems, requiring fast decisions that balance latency and energy consumption. Advances in Vehicular Edge Computing (\acrotip{VEC}) architectures have moved processing capacity closer to users, but high mobility and varying network loads make static heuristics inefficient. Consequently, recent literature has transitioned from traditional mathematical methods to Deep Reinforcement Learning (\acrotip{DRL}) approaches, seeking allocation policies that adapt dynamically to the environment.

The work presented in~\cite{vieira2025foff} developed the \acrotip{FOFF} framework based on Multi-Criteria Decision Making (\acrotip{MCDM}). The solution applied the Fuzzy-TOPSIS technique to rank edge servers (\acrotip{VEC}) and neighboring vehicles based on response time and energy consumption. Simulations demonstrated that the approach achieved significant energy savings without the need for a prior training phase. However, models based exclusively on logical rules or heuristic evaluations fail in non-stationary vehicular scenarios, as they cannot anticipate fluctuations in queue dynamics or learn from performance history.

The study conducted by~\cite{tan2018} explored machine learning to handle user mobility in vehicular networks. The research proposed a framework based on Deep Q-Networks (\acrotip{DQN}) to decide allocation at roadside units and manage task offloading. The results indicated an effective reduction in communication and computation costs. The main limitation of this approach lies in its centralized nature, requiring simultaneous control of all vehicles in the network and causing an exponential increase in computational complexity as vehicular density grows.

The work by~\cite{pang2024} developed a multi-task offloading algorithm using the Dueling Double Deep Q-Network (\acrotip{D3QN}) architecture. The method focused on minimizing response time and energy consumption in highly requested edge environments. Evaluations showed that the model reduced system latency more effectively than traditional random approaches. However, being based on strictly discrete action spaces, the algorithm faces scalability difficulties when applied to continuous resource allocation problems, limiting the precision of computational power distribution.

Research presented by~\cite{fofana2024} proposed an offloading solution based on a four-stage Stackelberg game integrated with a reinforcement learning model. The strategy evaluates vehicle state in real-time to optimize offloading time and expense in a coordinated manner. Experiments confirmed that the solution achieved balance in system utility and improved task success probability. Despite the gains, reliance on complex negotiation processes between multiple agents introduces signaling overhead that degrades performance in dense traffic and intermittent communication scenarios.

The study by~\cite{shi2023} sought to solve the problem of continuous resource allocation by applying the Deep Deterministic Policy Gradient (\acrotip{DDPG}) algorithm. The approach used an actor-critic structure to reduce overall system cost in a collaborative cloud and vehicular edge environment. The algorithm converged quickly and reduced operational costs compared to value-only methods. The failure of this modeling is the lack of consideration for structural dependencies between tasks and the use of zero-padding in state observations to handle a dynamic number of vehicles, forcing the neural network to process spurious artifacts and noise.

The work developed by~\cite{gao2022} implemented an adaptive scheme using multi-agent reinforcement learning in vehicular fog computing. The solution distributed decision-making directly among vehicles to reduce latency and energy consumption in a decentralized manner. Simulations demonstrated superiority in local adaptation capacity compared to single-agent guided architectures. The critical limitation of this architecture is the synchronization challenge and the difficulty of ensuring stable convergence when multiple vehicles change the environment simultaneously, generating instability in the global routing policy.

The study conducted by~\cite{zheng2023} integrated digital twin technology with vehicle networks to enable proactive artificial intelligence-assisted offloading. The work employed the Asynchronous Advantage Actor-Critic (\acrotip{A3C}) algorithm to optimize decisions based on a continuous virtual replica of the physical vehicular environment. The research reported rapid network convergence and considerable reduction in system cost. The fundamental restriction of this model is the extremely high computational overhead required to keep the digital twin synchronized in real-time, a factor that frequently makes its smooth execution on mobile devices unfeasible.

Research by~\cite{nie2023} directed reinforcement learning toward computer vision tasks and proposed an algorithm focused on object detection in autonomous vehicles. The model combined \acrotip{DRL} with piecewise linearization to specifically maximize a utility function based on execution time. The results pointed to significant improvements in response time for critical image processing. However, the modeling was built monolithically for a specific type of application, failing to consider the priorities of different competing tasks and neglecting congestion issues caused by vehicular density.

The study presented by~\cite{liu2024} focused on offloading complex tasks modeled as directed acyclic graphs using the Soft Actor-Critic (\acrotip{SAC}) algorithm. The framework sought to actively balance delay and energy expenditure for applications with high structural dependency between their subcomponents. Evaluation attested that stochastic entropy maximization avoided convergence to local optima and improved execution utility. The main restriction of this model for the automotive context is the premise of a perfectly static environment, completely ignoring disconnection rates and communication channel variability.

A hybrid approach was developed by~\cite{wu2024}, which proposed the \acrotip{TOLB} scheme for offloading and load balancing by integrating Twin Delayed Deep Deterministic Policy Gradient (\acrotip{TD3}) with the Technique for Order of Preference by Similarity to Ideal Solution (\acrotip{TOPSIS}). The methodology used the learning algorithm for partitioning proportion decisions and multi-criteria calculation to select the target server. Experiments proved the reduction of processing delay and decrease in load imbalance in the edge infrastructure. The limitation of this approach lies in using multi-criteria filtering only as a tactical post-processing element, keeping the observation space raw and rigid at the agent's input and hindering spatial generalization in dense networks.

The proposed \acrotip{SAC}-\acrotip{TOPSIS} architecture in this intermediate stage differs substantially from these approaches, as it consolidates the \acrotip{TOPSIS} method internally within the observation edge, acting strictly as a state engineering module. This module evaluates the edge server (\acrotip{VEC}) and neighboring vehicles based on Job Completion Time (\acrotip{JCT}), energy consumption, and distance to the computing node to identify the best offloading candidate. The computed output condenses these characteristics into a strictly compact continuous state encoding as a 6D vector that is perfectly scale-invariant, eliminating the dependence on zero-padding. The experimental evaluation compared five \acrotip{DRL} families (\acrotip{SAC}, Proximal Policy Optimization (\acrotip{PPO}), \acrotip{TD3}, \acrotip{DQN}, and Advantage Actor-Critic (\acrotip{A2C})) and seven baseline heuristics (\acrotip{VEC}, Random, Round-Robin, Gini, Jain, Greedy-Fair, and Greedy-Unfair), confirming that pairing this structured space yields consistent gains. This validates the central thesis that compositional state engineering developed via \acrotip{TOPSIS} acts as the primary driver for the generalization capability exhibited by the intelligent agent in variable traffic environments. In summary, the main aspects of related works are presented in Table~\ref{tab:related_works}.

\begin{table*}[tpb]
\centering
\caption{Comparison of Offloading Solutions -- A comparative evaluation of offloading strategies proposed in previous studies.}
\label{tab:related_works}
\renewcommand{\arraystretch}{1.3}
\begin{tabularx}{\textwidth}{c X X X X}
\hline
Ref. & Strategy & Communication Technology & Achievement & Limitation \\
\hline
\cite{vieira2025foff} & Fuzzy-\acrotip{TOPSIS} & \acrotip{VEC} & Significant energy savings without prior training phase & Fails to adapt to non-stationary vehicular scenarios \\
\cite{tan2018} & \acrotip{DQN} & Vehicular Networks & Effective reduction in communication and computation costs & Exponential increase in complexity due to centralized nature \\
\cite{pang2024} & \acrotip{D3QN} & \acrotip{IoV} and \acrotip{MEC} & More effective system latency reduction & Scalability difficulties in continuous allocation problems \\
\cite{fofana2024} & \acrotip{DRL} + Stackelberg Game & Vehicular Networks & System utility and task success probability balance & High signaling overhead in dense traffic scenarios \\
\cite{shi2023} & \acrotip{DDPG} & Vehicular Cloud and Edge & Fast convergence and operational cost reduction & Use of zero-padding and noise processing \\
\cite{gao2022} & \acrotip{MARL} & Vehicular Fog & Superiority in decentralized local adaptation capacity & Synchronization challenge and unstable global network convergence \\
\cite{zheng2023} & \acrotip{A3C} + Digital Twins & \acrotip{IoV} & Fast network convergence and system cost decrease & Extremely high computational overhead for synchronization \\
\cite{nie2023} & \acrotip{DRL} + Piecewise linearization & Autonomous Vehicles & Significant response time improvements for image processing & Neglects congestion issues and task priorities \\
\cite{liu2024} & \acrotip{SAC} & Vehicular Cloud & Avoids local optima convergence and improves utility & Static environment premise and channel variability ignorance \\
\cite{wu2024} & \acrotip{TD3} + \acrotip{TOPSIS} & \acrotip{VEC} & Processing delay and load imbalance reduction & Multi-criteria filtering used only as tactical post-processing \\
\hline
\end{tabularx}
\end{table*}
